\documentclass[10pt,a4paper]{amsart}
\usepackage[margin=19mm]{geometry}
\usepackage{amsmath,amssymb,mathtools}
\usepackage[scr=rsfs,cal=euler]{mathalfa}
\usepackage{xfrac}
\usepackage[unicode,hidelinks,bookmarks=false]{hyperref}
\newcommand{\ie}{i.e.\ }
\newcommand{\cf}{cf.\ }
\newcommand{\resp}{respectively\ }
\newcommand{\Subsection}{\subsection}
\providecommand{\nobreakdash}{\nobreak\hbox{-}}
\setlength{\emergencystretch}{2em}
\multlinegap0pt
\usepackage{bm}
\usepackage{bbm}
\usepackage{dsfont}
\usepackage{xspace}
\usepackage{cancel}
\usepackage{mathtools}
\usepackage{amsthm}
\makeatletter
\@ifundefined{lemma}{\newtheorem{lemma}{Lemma}}{}
\makeatother
\title[Convex bodies with pairs of sections associated by reflectio]{Convex bodies with pairs of sections associated by reflections}
\author{Efren Morales-Amaya}
\date{Source: arXiv:2407.02755v2 (CC BY 4.0)}
\begin{document}
\maketitle

\noindent\emph{Source and licence.} Convex bodies with pairs of sections associated by reflections. Version arXiv:2407.02755v2 (CC BY 4.0), distributed under the Creative Commons Attribution 4.0 International licence, as recorded in the arXiv licence metadata for this exact version and shown on the abstract page https://arxiv.org/abs/2407.02755v2. Full rights notice: licenses/arxiv-2407.02755-CC-BY-4.0.txt.\par\smallskip

\noindent\textit{Editorial scope.} Counted unit: the Lemma whose source label is \texttt{orto} in the article (source0.tex lines 210--212, source label \texttt{orto}), delivered together with the article's own complete original proof of it (source0.tex lines 214--228, one \texttt{proof} environment of 15 lines). No mathematics is written, repaired or re-derived: statement and proof are the article's own text line for line. Required context retained uncounted: none. Every definition, display and label of those lines is retained unchanged. Omitted from this package and not redistributed: 1--209, 213--213, 229--556. Nothing inside a retained excerpt is omitted or altered: every retained body is delivered exactly as the article's lines are written, so no omission receipt of any kind accompanies this package. Citations: the retained excerpts carry no citation command at all, so no bibliography entry of the article is transcribed and no reference is redistributed. Reference adaptation: none was needed and none was made; every reference inside the delivered unit resolves inside it, so no unresolved reference or citation is produced. Printed slips kept literal: no printed word, symbol, hypothesis or number of the counted statement or of its proof was altered. Layout and class: the article's own class is \texttt{article} with packages including amsmath, amssymb, amsthm, geometry, enumitem, graphicx, graphics, hyperref, amsfonts; the wrapper is the lane's pinned \texttt{amsart} editorial wrapper at 10pt on A4, which restates the theorem-like environments the retained text needs and adds the article's own macro definitions that the retained text needs (none), with extra wrapper packages (bm, bbm, dsfont, xspace, cancel, mathtools). The delivered numbering is the unit's own. Assets: no retained span contains a figure environment, a graphics-inclusion command, a PDF-page inclusion, a table or a TikZ picture; no image file is copied into this package, the package declares no asset dependency and no image file is redistributed with it. Licence: version arXiv:2407.02755v2 of this article is distributed under the Creative Commons Attribution 4.0 International licence, as recorded in the arXiv licence metadata for this exact version and shown on the abstract page https://arxiv.org/abs/2407.02755v2. A full rights notice is delivered at licenses/arxiv-2407.02755-CC-BY-4.0.txt. Characters: every retained excerpt is pure ASCII, so the delivered engine, XeLaTeX with its default Latin Modern text font, reports no missing character.
	\begin{lemma}\label{orto}
		If $p_1 \notin K_1$, then $p_2 \notin K_2$ and the line $L$ is perpendicular to $H$.
	\end{lemma}

	\begin{proof}
		We first claim that $p_2 \notin \operatorname{int}(K_2)$. Let $G_1$ be the projection of $K_1$ from $p_1$ onto $H$. Let $M$ be a supporting line of $G_1$ touching it at a unique point. Then $\pi_1(M) \cap K_1$ is at most 1-dimensional. If $p_2 \in \operatorname{int}(K_2)$, then $\pi_2(M) \cap K_2$ would be 2-dimensional, contradicting the existence of a reflection mapping $\pi_2(M) \cap K_2$ onto $\pi_1(M) \cap K_1$. Thus, $p_2 \notin \operatorname{int}(K_2)$.
		
		Let $G_2$ be the projection of $K_2$ from $p_2$ onto $H$. We show $G_1 = G_2$. Assume otherwise, say $x \in G_1 \setminus G_2$. Without loss of generality, let $x \in \operatorname{int}(G_1)$. Let $W \subset H$ be a line separating $x$ from $G_2$ and let $M$ be a line parallel to $W$ through $x$. Since $x \in \operatorname{int}(G_1)$, $\pi_1(M) \cap K_1$ contains interior points. But $M \cap G_2 = \emptyset$, so $\pi_2(M) \cap K_2 = \emptyset$, a contradiction. Hence, $G_1 = G_2$.
		
		Suppose $L$ is not perpendicular to $H$. Let $\alpha, \beta$ be the projections of $p_1, p_2$ onto $H$, respectively. Since $L$ is not perpendicular to $H$, $\alpha \neq \beta$.
		
		We claim there exists a point $A \in \partial G_1$ and two perpendicular lines $M_A, N_A$ such that $M_A$ is a supporting line to $G_1$ with $M_A \cap \partial G_1 = \{A\}$, $N_A$ passes through $A$, and $N_A \cap \operatorname{int}[\alpha, \beta] \neq \emptyset$.
		
		Let $w \in \operatorname{int}[\alpha, \beta]$. Choose a disc $B \subset \mathbb{R}^2$ centered at $w$ such that $G_1 \subset \operatorname{int}(B)$ (this is possible since $G_1$ is compact). Shrink the radius of $B$ until $B \cap \partial G_1 \neq \emptyset$. Let $A \in B \cap \partial G_1$, and let $M_A$ be a supporting line of $B$ through $A$. Since $G_1 \subset B$, $M_A$ is also a supporting line of $G_1$. The tangency condition implies $M_A \cap \partial G_1 = \{A\}$.
		
		Let $N_A := L(A, w)$, and let $\Delta$ be the plane perpendicular to $H$ containing $N_A$, with open half-spaces $\Delta^+, \Delta^-$ such that $p_1 \in \Delta^+$ and $p_2 \in \Delta^-$. Define $\Gamma_1 := \operatorname{aff}\{M_A, p_1\}$, $\Gamma_2 := \operatorname{aff}\{M_A, p_2\}$, and let $S_M$ be the reflection with $S_M(\Gamma_2 \cap K_2) = \Gamma_1 \cap K_1$.
		
		Since $L(p_2, A)$ is a supporting line of $K_2$, it intersects $K_2$ at some $\tau \in \partial K_2 \cap \Delta^-$. Then $S_M(\tau) \in \partial K_1 \cap \Gamma_1$. Thus, $L(p_1, S_M(\tau))$ is a supporting line of $K_1$, but it cannot pass through $A$, contradicting the choice of $M_A$. Therefore, $L$ must be perpendicular to $H$.
	\end{proof}

\end{document}
